\begin{apartats}

\apartat[1,25]{0,75}
Determina el domini de les funcions següents:
\begin{graella}{3}
  \sa f(x)=\frac{x+1}{x^2-4} & \sa g(x)=\sqrt{x^2-x-6} & \sa h(x)=\ln(3-x)
\end{graella}

\begin{solucio}
i) Cal que el denominador no s'anul·li: $x^2-4=0$ en $x=\pm2$. Domini:
$\mathbb{R}\setminus\{-2,2\}$.\\
ii) Cal $x^2-x-6=(x-3)(x+2)\ge0$: domini $(-\infty,-2]\cup[3,+\infty)$.\\
iii) Cal $3-x>0$: domini $(-\infty,3)$.
\end{solucio}

\apartat[1,25]{1}
Determina el domini i els punts de tall amb els eixos de
\[
f(x)=\frac{x^2-9}{x+2} .
\]

\begin{solucio}
\textbf{Domini}: $x+2\neq0$, és a dir $\mathbb{R}\setminus\{-2\}$.\\
\textbf{Tall amb l'eix $OY$}: $f(0)=\dfrac{-9}{2}$, el punt $\left(0,-\tfrac92\right)$.\\
\textbf{Talls amb l'eix $OX$}: $f(x)=0$ quan el numerador s'anul·la i el denominador no:
$x^2-9=0$ dona $x=3$ i $x=-3$. Els punts són $(3,0)$ i $(-3,0)$.
\end{solucio}

\begin{nomesllarg}
\apartat{0,75}
Determina el domini de
\[
k(x)=\frac{\sqrt{x+1}}{x-2} .
\]

\begin{solucio}
Cal complir dues condicions alhora: $x+1\ge0$, és a dir $x\ge-1$, i $x-2\neq0$.\\
Domini: $[-1,+\infty)\setminus\{2\}$.
\end{solucio}
\end{nomesllarg}

\end{apartats}
